\chapter{网络流量智能应用分析系统设计与实现}[Design and Implementation of the Intelligent Network Traffic Application Analysis System]

本章在第3章需求分析的基础上，给出本系统的总体架构与模块化实现方案，并结合关键算法逐步说明各组成部分的设计取舍与运行机制。系统以离线PCAP文件为统一输入，最终输出包含已知流量标签、未知流量复合标签、应用与服务类型聚合结果以及可视化报告的端到端审计成果。

\section{系统总体架构设计}[Overall System Architecture Design]

系统总体上采用“两阶段混合\,+\,插件式扩展”的架构思想。第一阶段为确定性预处理与已知流量过滤，第二阶段为面向未知流量的模型辅助识别与智能体语义分析；在两阶段之间，预训练分类模型以可插拔形式接入，作为面向未知流量的统计层证据来源。系统在工程上被划分为预处理过滤、流量分类、智能分析、结果合并与可视化四个核心模块，并由统一的流水线驱动器串联，整体逻辑结构如图\ref{fig:ch4-architecture}所示。

\begin{figure}[h]
	\centering
	\fbox{\parbox[c][6cm][c]{0.8\textwidth}{\centering 占位图：系统总体架构。包含原始PCAP输入、预处理过滤、未知流PCAP切分、分类模型推理、智能体推断、结果合并与可视化六大节点，节点之间通过结构化数据契约相互衔接。}}
	\bicaption[ch4-architecture]{}{系统总体架构}{Fig.}{Overall architecture of the system}
\end{figure}

各模块之间通过严格定义的JSON数据契约进行交互，主要包括预处理结果契约、智能体输入契约、智能体输出契约和最终报告契约。不同模块在可独立运行的基础上，仅通过约定的接口对象进行耦合，从而保证系统具有良好的可调试性、可替换性与可扩展性。模块职责边界划分如下。

预处理过滤模块只承担规则化、可复现的工作，包括协议解析、流聚合、元数据抽取、白名单匹配、未知流PCAP切分以及分类模型回填，不参与任何语义层面的推断。流量分类模型模块以独立子项目形式提供，对外暴露统一的适配函数接口，支持在不修改主流程的前提下替换或停用模型。智能分析模块基于本课题集成的智能体框架，对预处理无法确定的流量进行多源证据融合与语义判断，输出结构化复合标签。结果合并与可视化模块不参与判别，只负责把已知与未知结果按流唯一键聚合，并以图形化方式呈现给审计人员。

为了在不同实验环境下保持稳定运行，系统支持降级执行：当分类模型权重缺失或外部模型服务不可用时，预处理结果仍可独立输出已知与未知划分；当智能分析模块所依赖的大语言模型不可用时，流水线会跳过Agent阶段并标记为已跳过状态，最终报告仍能完整生成。这种弱耦合的工程结构使得系统既能在算力与外部资源充足时执行完整推理链，也能在受限环境下提供可用的基础审计能力。

\section{已知流量过滤模块设计与实现}[Design and Implementation of the Known Traffic Filtering Module]

已知流量过滤模块是连接原始流量与后续智能分析的入口环节。其设计目标在于：以确定性、低延迟、可解释的方式过滤大部分常见流量，并在不丢失审计上下文的前提下，把无法可靠判定的流量整理为后续模块的标准化输入。

模块内部采用流水线式处理结构，其总体流程如图\ref{fig:ch4-preprocess}所示。预处理调度器依次完成日志生成、流聚合、元数据增强、规则判别、未知流切分与分类模型回填等步骤。具体流程包括：以原始PCAP为输入运行Zeek，得到连接、SSL、HTTP、DNS、X.509等结构化日志；调用nDPI生成包含五元组、协议栈、SNI、风险标志和流统计字段的流级CSV；随后以nDPI流为基本单元，将Zeek日志中携带相同连接标识或同源五元组的记录合并到同一条流上，构造出统一的元数据视图。

\begin{figure}[h]
	\centering
	\fbox{\parbox[c][6cm][c]{0.8\textwidth}{\centering 占位图：已知流量过滤模块流程。从原始PCAP出发，依次经过Zeek与nDPI解析、流聚合与元数据增强、判别规则、已知/未知划分、未知流PCAP切分与分类模型回填，最终输出结构化预处理结果。}}
	\bicaption[ch4-preprocess]{}{已知流量过滤模块处理流程}{Fig.}{Processing flow of the known traffic filtering module}
\end{figure}

对Zeek日志的对齐采用双键索引策略。一方面以连接唯一标识为键将SSL、HTTP、DNS等记录归并到对应连接；另一方面分别以严格五元组键和方向无关的规范化五元组键建立连接索引，并在多候选时按时间戳就近匹配，避免因双向流方向反转或源端口随机化导致的误关联。该设计使得Zeek与nDPI两套异构输出能够稳定地对应到同一条网络流。

为了避免后续模块依赖松散字段，模块基于Pydantic定义了统一的流元数据模型，并嵌套HTTP、DNS、TLS、流统计、PCAP切分状态以及分类模型结果等子结构，覆盖五元组、传输层协议、加密状态、应用层语义字段、证书链指纹、流量统计特征和未知流PCAP产物等约30余项可观测信息。在解析阶段，模块按字段类型完成空值清洗、域名归一化、向量解析与布尔值兼容处理，确保进入判别逻辑的数据结构稳定一致。

判别逻辑按从强到弱的优先级依次执行多条规则，整体过程如算法\ref{alg:ch4-classify}所示。命中已知规则的流量被赋予较高的预处理置信度，未知流量则赋予较低的基线置信度，作为后续融合的初值。SNI白名单支持精确匹配与通配符匹配，命中条目可直接生成形如 $\langle \mathrm{app} \rangle\!:\!\langle \mathrm{service\_type} \rangle$ 的复合标签；系统协议集合涵盖DNS、DHCP、ARP、ICMP、NTP、SSDP、mDNS等常见管理类协议；高置信应用提示则面向nDPI已能可靠识别的具体应用进行直接映射。



\begin{algorithm}[h]
	\caption{已知流量过滤判别流程}\label{alg:ch4-classify}
	\KwIn{流元数据 $f$, 已知SNI白名单 $\mathcal{W}$, 系统协议表 $\mathcal{S}$, 高置信应用表 $\mathcal{H}$}
	\KwOut{三元组 $(is\_known, label, reason)$}
	
	\If{$f.sni \in \mathcal{W}$}{
		\Return{$(\mathrm{true}, \text{lookup}(\mathcal{W}, f.sni), \text{SNI命中白名单})$}\;
	}
	\If{$f.ndpi\_app \in \mathcal{S}$}{
		\Return{$(\mathrm{true}, f.ndpi\_app, \text{系统协议快速过滤})$}\;
	}
	\If{$f.ndpi\_app \in \mathcal{H}$}{
		\Return{$(\mathrm{true}, f.ndpi\_app, \text{高置信应用映射})$}\;
	}
	\If{$f$ 属于广播、组播或链路本地范围}{
		\Return{$(\mathrm{true}, \text{system:local\_discovery}, \text{链路本地噪声})$}\;
	}
	\If{$f.http\_host \in \mathcal{W}$}{
		\Return{$(\mathrm{true}, \text{lookup}(\mathcal{W}, f.http\_host), \text{HTTP Host命中白名单})$}\;
	}
	\uIf{$\neg f.is\_encrypted$}{
		\Return{$(\mathrm{false}, \mathrm{nil}, \text{明文流量待深度解析})$}\;
	}
	\uElseIf{$f.sni = \varnothing$}{
		\Return{$(\mathrm{false}, \mathrm{nil}, \text{加密流量缺失SNI})$}\;
	}
	\Else{
		\Return{$(\mathrm{false}, \mathrm{nil}, \text{加密流量SNI未命中白名单})$}\;
	}
\end{algorithm}

模块同时为每次任务生成稳定的任务标识，并以该标识为命名空间在结果目录下完成未知流PCAP切分。这一过程通过无第三方依赖的方式解析经典PCAP头部、按字节序与时间戳精度遍历包记录，再以五元组流键与时间区间联合判定每个数据包归属，最终生成单流PCAP文件以及描述全部切分结果的任务清单。任务清单将用于流量分类模型的批量推理输入，并在元数据层面保留切分路径与切分状态字段，以便审计人员追溯原始报文。

模块最终将处理结果以三份JSON落盘：完整的预处理结果用于审计与下游消费，已知子集与未知子集分别保存便于按需读取。整个过滤过程对处理对象做了显式裁剪：通过字段精简函数移除Zeek原始记录、X.509裸证书等大块字段，避免将冗余信息传递给智能分析阶段，从而在控制下游令牌消耗的同时保留审计可追溯性。

\section{流量分类模型模块设计与实现}[Design and Implementation of the Traffic Classification Model Module]

流量分类模型模块负责在预处理阶段为未知流量提供基于深度学习的服务类型证据。模块在工程层面被设计为可插拔组件，主流程通过统一的适配函数接口与底层模型解耦，使得模型既可以独立训练与演化，也可以在不修改主程序的前提下替换或关闭。

\subsection{适配层与模型架构}

适配层提供统一的入口函数，接收预处理阶段产出的任务清单为输入，并根据系统配置中的启用标志决定是否触发模型推理；当启用时，按配置项动态加载形如“模块名:函数名”的具体适配函数并调用之。适配函数的返回值经过统一化处理后，按流序号与流唯一键两种键形成字典索引，再回填到对应未知流的分类模型结果字段中。当模型不可用、未配置或运行失败时，系统会返回带有错误信息的占位结果，使后续模块仍能正常运行。这种以协议而非耦合代码的方式接入模型，符合需求分析中提出的可扩展性与容错性要求。

模型整体采用预训练Transformer编码器结构，其总体架构如图\ref{fig:ch4-model}所示。模型以原始报文字节序列为输入，按传输方向将一条会话流划分为若干连续报文窗口；每个窗口截取若干连续报文的有效载荷并以双字节为单位进行十六进制分词，再附加分类、分隔与填充等特殊标记构造成定长令牌序列。模型主体由12层Transformer编码块构成，隐藏维度为768，每层配备12个注意力头，自注意力机制可同时建模窗口内字节级依赖与窗口间通信时序依赖。模型权重在大规模无标签加密流量上完成自监督预训练，再针对本课题的服务类型分类任务进行下游微调。最终输出五类服务类型概率分布，分别为大批量传输（bulk-transfer）、交互控制（interactive）、流媒体（stream）、隧道代理（vpn）和一般网页访问（web）。这一类别空间与智能分析模块所使用的服务类型词表保持完全一致，便于后续证据融合。

\begin{figure}[h]
	\centering
	\fbox{\parbox[c][5cm][c]{0.8\textwidth}{\centering 占位图：流量分类模型架构。原始报文经分词层与嵌入层后输入12层Transformer编码块，最终经分类头输出五类服务类型概率。}}
	\bicaption[ch4-model]{}{流量分类模型架构}{Fig.}{Architecture of the traffic classification model}
\end{figure}

\subsection{模型微调过程}

预训练模型虽具备较强的通用流量表示能力，但其类别空间与本系统所需的服务类型层标签并不完全一致，因此需要在加载预训练权重的基础上完成下游微调，使模型从通用表示平稳迁移到具体分类任务。

微调阶段采用全参数更新策略，以充分释放预训练表示在加密流量任务上的潜力。训练数据由公开数据集与自采集流量混合构成，覆盖多种加密协议、运行平台与正常/异常流量类型，具体来源如表\ref{tab:ch4-finetune-dataset}所示。为缓解类别不平衡，每类样本设置5000条报文的上限采样策略，并对模糊语义混合样本进行剔除，以降低标签噪声带来的拟合偏差。每个训练样本由若干连续报文构成的窗口序列组成，并经分词后限制最大长度为512个令牌。

\begin{table}[h]
	\centering
	\caption{微调阶段使用的数据来源}\label{tab:ch4-finetune-dataset}
	\small
	\begin{tabular}{ll}
		\toprule
		数据集 & 主要内容 \\
		\midrule
		ISCXVPN2016 & 隧道封装与直连流量，覆盖多类应用服务 \\
		ABNT Dataset & 桌面端与移动端常见应用流量 \\
		CSTNET-TLS1.3-Pcaps & 基于TLS\,1.3协议的网页与应用流量 \\
		CICIDS2017 & 公开网络入侵检测流量 \\
		USTC-TFC2016 & 正常应用与恶意攻击流量 \\
		自采集校园网流量 & 真实环境下的多类业务流量 \\
		\bottomrule
	\end{tabular}
\end{table}

微调采用AdamW优化器，初始学习率为$2 \times 10^{-5}$，批量大小为32，最大训练轮数为10。为提升泛化能力并控制过拟合风险，分类头全连接层引入0.1的Dropout，训练过程中实时监控验证集损失与F1指标并启用早停机制。具体超参数设置如表\ref{tab:ch4-finetune-hparams}所示。微调过程整体流程见算法\ref{alg:ch4-finetune}。

\begin{table}[h]
	\centering
	\caption{微调阶段关键超参数}\label{tab:ch4-finetune-hparams}
	\small
	\begin{tabular}{ll}
		\toprule
		参数 & 设置值 \\
		\midrule
		优化器 & AdamW \\
		初始学习率 & $2 \times 10^{-5}$ \\
		批量大小 & 32 \\
		最大训练轮数 & 10 \\
		最大输入序列长度 & 512 \\
		分类头Dropout比例 & 0.1 \\
		每类采样上限 & 5000条报文 \\
		权重微调策略 & 全参数微调 \\
		早停监控指标 & 验证集损失与F1 \\
		\bottomrule
	\end{tabular}
\end{table}

\begin{algorithm}[h]
	\caption{流量分类模型微调流程}\label{alg:ch4-finetune}
	\KwIn{训练集 $\mathcal{D}_{tr}$, 验证集 $\mathcal{D}_{val}$, 预训练权重 $\theta_0$, 最大轮数 $E$}
	\KwOut{微调后权重 $\theta^*$}
	
	$\theta \gets \theta_0$；初始化最佳指标 $m^* \gets 0$, 早停计数 $k \gets 0$\;
	\For{$e = 1,2,\dots,E$}{
		\For{每个批次 $\mathcal{B} \subset \mathcal{D}_{tr}$}{
			构造连续报文窗口并完成分词与填充\;
			前向计算交叉熵损失 $L$\;
			通过 AdamW 更新参数 $\theta$\;
		}
		在 $\mathcal{D}_{val}$ 上评估宏平均 F1，记为 $m_e$\;
		\eIf{$m_e > m^*$}{
			$m^* \gets m_e$, $\theta^* \gets \theta$, $k \gets 0$\;
		}{
			$k \gets k + 1$\;
		}
		\If{$k$ 达到早停阈值}{
			\textbf{break}\;
		}
	}
	\Return{$\theta^*$}\;
\end{algorithm}


为提升模型在域外流量上的稳定性，本系统在训练数据准备阶段进行了两轮迭代。首轮训练直接使用原始数据集，模型在闭集验证上表现出极高的拟合速度，但在域外样本上出现明显的分类偏置，提示模型可能记忆了数据集特定的统计噪声而非学习到通用语义。次轮训练通过剔除语义混合样本、对每类样本进行上限采样并引入早停机制，使模型从“样本记忆”转向“特征泛化”，在不显著降低闭集准确率的前提下显著缓解了高置信度误判现象，从而获得更稳定的下游证据。

\subsection{推理与回填}

适配函数的具体执行流程分为两个阶段。第一阶段为推理数据生成：按配置的载荷长度、窗口报文数与样本上限将单流PCAP批量转换为推理用的TSV文件以及记录每条样本对应原始PCAP的清单。第二阶段为前向推理：加载微调后的模型权重与词表，按预设批大小执行推理，并将每个未知流的预测标签、概率和Top-K候选写入按流组织的JSON文件；最后由对齐函数按流序号与PCAP路径两种方式回填到任务清单，缺失预测的流则填充“无预测”状态。

需要强调的是，分类模型在系统中并不作为最终判定依据，而是作为下游智能分析模块工具链中的一项低权重证据。这是因为模型在域外流量上仍存在置信度偏置和分布漂移问题，单独依赖模型容易出现高置信误分类现象。本系统通过将模型输出限定为五类服务类型概率，并在Agent推理链中将其优先级置于SNI、主动获取与证书/IP证据之后，使模型既能在SNI不可见或元数据稀缺的场景下补充信息，又不会主导最终决策，从而实现确定性证据与统计性证据的互补。

\section{LLM Agent 智能分析模块设计与实现}[Design and Implementation of the LLM Agent Intelligent Analysis Module]

智能分析模块将预处理阶段保留的未知流交给基于大语言模型的智能体，通过工具调用与多源证据融合生成形如 $\langle \mathrm{app} \rangle\!:\!\langle \mathrm{service\_type} \rangle$ 的复合标签，并附带置信度与证据摘要，是系统中唯一承担语义层推断的模块。

为兼顾可维护性与隔离性，模块以适配层方式封装一个独立的智能体引擎子项目，并通过异步运行接口提供统一入口。运行时，主程序仅向引擎传入工作区路径与配置文件，所有提示词构造、会话管理、工具注册等逻辑均由引擎负责。模块同时为智能体准备了独立的工作目录，其中包含角色契约、工具优先级说明、风格约束、项目背景、技能定义以及引擎配置。引擎在运行过程中读取上述文档作为上下文，并在隔离工作区下完成会话与缓存写入，其文件操作通过工作区写入约束选项被严格限制，避免对系统其他目录造成意外修改。

为控制单次推理的上下文规模，输入构造模块按设定的分块上限将未知流列表切分写入工作区输入目录。每条流在落盘前进行字段裁剪，仅保留智能体推断所需的关键字段：五元组、加密状态、SNI、HTTP相关字段、DNS查询、TLS版本与SAN域名、协议栈识别结果、流统计特征以及分类模型的标签与概率。所有输入文件遵循统一的契约，并附带服务类型词表与证据优先级链，作为引擎判断时的硬约束。智能体可用的服务类型词表如表\ref{tab:ch4-vocab}所示，与分类模型类别空间完全对齐。

在进入分块构造之前，模块还设置了一道可识别性预筛分。其设计依据在于：当一条流同时缺失SNI、HTTP主机、DNS查询且分类模型未给出有效输出、有效字节数低于阈值时，可观测特征已不携带任何能够支撑应用归属判断的语义信号，将其送入大语言模型的期望信息增益接近于零。基于这一观察，模块按"具备至少一类识别信号"为条件把未知流划分为可推断子集与无信号子集：前者按既定分块策略进入推理流程，后者直接生成附带具体缺失原因的空决策（如同时记录"缺少SNI/HTTP主机/DNS查询、分类模型无输出、有效字节数过低"等条目），不再经过模型调用。该预筛分既显著降低了批次输入规模与令牌消耗，也使下游审计能够明确区分"模型未能判定"与"系统判定无可用证据"两种本质不同的失败模式。

\begin{table}[h]
	\centering
	\caption{服务类型词表}\label{tab:ch4-vocab}
	\small
	\begin{tabular}{ll}
		\toprule
		标签 & 含义 \\
		\midrule
		bulk-transfer & 大批量数据传输，如文件下载、备份、CDN大对象 \\
		interactive & 交互或控制类，如聊天、远程登录、即时通信信令 \\
		stream & 音视频流，如直播、点播、语音通话 \\
		vpn & 隧道或代理类流量 \\
		web & 一般网页或小流量HTTP/HTTPS访问 \\
		\bottomrule
	\end{tabular}
\end{table}

智能体的工具能力按层组织。基础读写层提供文件读写、通配符匹配、文本检索等操作，承担输入读取与决策落盘任务；信息获取层提供网页主动获取与搜索引擎调用，用于在SNI已知但未命中白名单时访问目标站点获取标题、响应头与证书摘要；域名情报层在引擎内部以原生工具形式集成了注册信息查询、多类型DNS记录解析、DNS健康度检查与综合域名分析等操作，分别支撑不同强度的证据获取需求；分类模型证据则以静态字段方式直接写入输入文件，不再独立暴露为工具，从而避免重复调用与上下文膨胀。所有工具均按统一的JSON Schema进行参数校验，运行时输出结构化结果以便后续解析。

在工具能力配置之外，模块进一步引入两类机制以保证工具被正确、节制地使用。第一类为工具能力曝光约束。本课题在前期实验中观察到，大语言模型即使在工具列表中已注册专用工具，仍倾向于反复调用最显眼的通用工具（如通用搜索），导致已注册的专业能力实际上处于失活状态。为此模块在角色契约、技能定义与提示词模板中以二维表的形式给出"任务-首选工具-禁用工具"对应关系，明确通用搜索类工具的禁用边界与专用情报工具的优先调用条件，并附带典型参数示例，使智能体即便不阅读完整技能文档也能在系统提示词中直接读到工具路由信息。第二类为运行时调用预算门控。模块在引擎运行时层面设置三道正交检查：相同工具与参数组合的重复调用上限、跨工具对同一目标的查询合计上限以及参数类型与工具语义的匹配检查（如对仅支持域名的工具拒绝直接传入IP字面量）。任何被拦截的调用均被替换为携带可读说明的合成错误结果回灌给模型，使其在同一推理轨迹内即时调整路径，避免重复试错带来的时延累积。两类机制分别从"曝光质量"与"调用上限"两个维度对工具使用行为进行约束，与既有的提示词层语义约束形成互补。

引擎的推理过程通过结构化提示词进行约束。提示词要求智能体先读取分块输入与角色文档，再按证据优先级逐条流推断，并以约定的结构写入指定输出路径。提示词同时给出明确的服务类型词表，并强调“证据不足时必须将复合标签置空，置信度归零”，以防止幻觉性输出。证据优先级链定义如下：
\begin{equation}\label{eq:ch4-priority}
	\mathrm{SNI} \succ \mathrm{active\_fetch} \succ \mathrm{cert/ip} \succ \mathrm{behavior\_stats} \succ \mathrm{classification\_model}
\end{equation}
式中，$\succ$ 表示前一证据来源在融合判断中的权重高于后者。该优先级既反映了不同证据的可靠性差异，也保证了模型预测在缺乏直接语义证据时仍能发挥补充作用。

智能体在每条流上的整体决策过程如算法\ref{alg:ch4-agent}所示。在分块执行层面，由于不同分块对应的推理任务之间互不依赖，运行驱动器将其建模为相互独立的异步协程，并通过信号量限定单进程下的最大并发度。这一设计利用了大语言模型调用属于I/O密集型这一特性，在不触发服务端速率限制的前提下并行处理多个分块；并发上限可通过环境变量进行配置，以适应不同提供方的服务约束。遇到单块失败时，运行驱动器将异常信息写入对应输出文件，从而保证整体批处理不会因为局部错误而中断。当多条流出现相同SNI或目标IP时，可由引擎在工作区缓存中直接复用结论，进一步降低重复调用成本。

\begin{algorithm}[h]
	\caption{智能体单条流量推断流程}\label{alg:ch4-agent}
	\KwIn{待标注流 $f$, 工具集 $\mathcal{T}$, 服务类型词表 $\mathcal{V}$}
	\KwOut{决策 $d = (\mathrm{app}, \mathrm{service\_type}, \mathrm{confidence}, \mathrm{evidence})$}
	
	初始化证据集合 $E \gets \varnothing$\;
	\uIf{$f.sni \neq \varnothing$ 且 SNI命中已知品牌}{
		将 SNI 证据加入 $E$, 设置较高置信度\;
	}
	\ElseIf{$f.sni \neq \varnothing$}{
		调用主动获取工具访问该域名，将标题与证书摘要加入 $E$\;
	}
	\If{仍缺乏强证据 且 $f.dst\_ip$ 公网可解析}{
		调用域名情报工具查询注册与 DNS 信息，加入 $E$\;
	}
	读取 $f.classification\_model$ 字段并以低权重加入 $E$\;
	以行为统计特征作为兜底证据加入 $E$\;
	按式\eqref{eq:ch4-priority}对 $E$ 进行加权融合\;
	\eIf{融合结果支持唯一应用与服务类型}{
		生成 $d$, 其中 $d.\mathrm{service\_type} \in \mathcal{V}$\;
	}{
		将 $d$ 各字段置空, $d.\mathrm{confidence} \gets 0$, 并在理由中记录证据缺失原因\;
	}
	\Return{$d$}\;
\end{algorithm}


智能体的输出由结果解析器进行校验与归一化。考虑到大语言模型生成的JSON常带有Markdown围栏、行内注释、悬挂逗号或多余前后缀等不规范成分，解析器采用渐进降级的解析链：优先尝试严格JSON解析，失败后依次执行围栏剥离、注释清除、悬挂逗号修复以及最外层花括号截取等修复操作，每一级修复均记录恢复说明以便后续审计追溯。解析得到的多种返回格式（裸列表、决策数组或单条字典）经统一展开后进入字段标准化阶段：拆分复合标签为应用与服务类型字段、把不在词表内的服务类型置空并清除复合标签、将置信度截断到$[0,1]$区间、规整证据列表的字段名。为防御大语言模型对输入流唯一键的幻觉性输出，解析器在调用方传入的合法流键集合上进行硬过滤，对不在该集合中的决策直接丢弃并写入错误列表，避免幻觉条目污染最终聚合结果。归一化后的对象再交给基于Pydantic的模式校验，校验失败时仅剔除越界字段而保留可用部分，避免因单条异常导致全量结果丢失。最终结果按统一契约合并去重，并附带运行过程中收集到的错误列表，供合并模块使用。

在结果解析完成后，模块还设置了一道证据兜底融合环节，用于弥合大语言模型与确定性分类器在置信度表达方式上的差异。大语言模型在证据稀疏时倾向给出空标签以避免幻觉，但流水线中已存在的预训练分类模型通常仍能给出具有一定置信度的概率分布。模块对所有最终标签为空的决策做二次扫描：若对应流的分类模型概率达到设定阈值，则按式\eqref{eq:ch4-fallback}合成兜底结果，应用字段固定为"unknown"以显式表达"协议层未知但行为类型可判"的语义；置信度由分类模型概率乘以衰减系数得到，以避免兜底证据在融合中被过度采信，并在证据链中保留分类模型来源标记，便于下游审计区分智能体直接判定与模型概率合成两类结果。
\begin{equation}\label{eq:ch4-fallback}
  \mathrm{conf}_{fb} = \alpha \cdot p_{cm}, \quad p_{cm} \geq \tau
\end{equation}
式中，$p_{cm}$ 为分类模型给出的最高概率，$\tau$ 为触发阈值，$\alpha \in (0,1)$ 为置信度衰减系数。该机制以低耦合方式提升了系统对证据稀疏未知流的整体覆盖率，同时未引入额外的大语言模型调用成本。

通过上述设计，智能分析模块实现了从未知流到可解释复合标签的端到端转换，同时以契约化输入输出、隔离工作区、严格工具约束与多重容错机制，将大语言模型固有的不稳定性控制在工程可接受的范围内。

\section{结果合并与可视化模块设计与实现}[Design and Implementation of the Result Merging and Visualization Module]

结果合并与可视化模块负责把预处理与智能分析两阶段的成果汇总为面向审计人员的报告与图表，是系统中直接产出可消费成果的环节。模块在职能上进一步拆分为三层：合并层负责把异构产物按流唯一键对齐为统一的最终报告契约，统计层负责由最终报告派生面向审计的多维度指标，呈现层则同时提供可独立分发的静态图表与可交互的Web仪表盘。

合并层以流唯一键为索引，将智能体的决策按流回填到对应的未知流对象上，并按"标签是否有效"与"是否携带具体原因"两个维度引入三分支处理。对于智能体给出有效复合标签的决策，应用、服务类型、置信度、理由、证据、工具调用线索等字段被原样写入未知流对象；对于智能体已对该流做出判断但因证据不足未能给出复合标签的决策（例如预筛分阶段标记的无信号流量与大语言模型主动声明证据不足的流量），合并层保留其特定的未标注原因与证据，仅将复合标签字段置空，避免具体诊断信息在合并阶段被丢失；对于既未出现在智能体输出中、也未通过预筛分生成占位决策的流量，合并层使用统一的占位模板填充以保持输出对象结构一致。这种三分支处理使得最终报告中的每一条未知流都携带可解释的状态信息，下游审计可清晰区分"已成功标注"、"已判定但证据不足"和"未进入推理流程"三类不同情形。合并完成后，模块在原有统计字段上追加智能体成功标注数量等概要指标，并将完整的最终报告以约定契约持久化。

统计层在最终报告上派生面向审计的扩展指标。对于每条流量记录，统计层从已知子集与未知子集合并的全集中提取流统计字段，计算总字节数、平均与最长持续时间，以及nDPI识别率（按非通用未知字符串过滤后判定的可识别比例）；对于聚合维度，按应用与服务类型分别统计计数及占总流数的百分比，并把未知流量按最终复合标签归并为可在表格中展示的分组结果。聚合规则简化形式如式\eqref{eq:ch4-aggregation}所示。
\begin{equation}\label{eq:ch4-aggregation}
  N_{a,s} = \sum_{i=1}^{N} \mathbb{1}[f_i.\mathrm{app} = a \,\wedge\, f_i.\mathrm{service\_type} = s]
\end{equation}
式中，$N$ 为流总数，$\mathbb{1}[\cdot]$ 为指示函数，$N_{a,s}$ 表示应用 $a$ 与服务类型 $s$ 的联合流数量。该量同时构成桑基图、应用-服务交叉表与未知结果分组表的基础。所有聚合结果在序列化时按出现频次降序排列，便于直接用于占比展示，并将相对百分比与绝对计数一同保留，使审计人员能够同时观察规模与构成。

呈现层提供"静态图表产物"与"交互式仪表盘"两条互补路径。静态图表产物基于Plotly构造三类离线HTML：以应用占比为输入生成空心环形饼图，反映应用层占比；以服务类型占比为输入生成柱状图，呈现服务类型词表上的分布；以应用与服务类型映射为输入生成桑基图，直观刻画应用与服务类型之间的归属关系。绘图后端在导入失败时返回空值而非抛出异常，使得在缺失绘图依赖的环境下整体流水线仍可继续；该路径默认关闭，仅在用户显式启用时生成，主要用于离线审计材料的归档与分发。交互式仪表盘则是默认呈现方式，由Web服务在每次任务完成后即时构建：前端首先以扁平结构展示由统计层派生的全局指标和按标签分组的未知结果表，并在浏览器侧通过同一图形库渲染应用占比图；当审计人员点击某条未知流时，前端发起单流下钻请求，由后端结合最终报告与本任务记录的输入分块、推理会话快照构造详情包，再以三栏视图呈现——概要栏展示该流的复合标签、置信度、理由、证据列表与工具调用线索；流特征栏展示五元组、TLS与DNS字段、字节统计以及分类模型的预判概率；推理过程栏按角色对智能体的多轮思考、工具调用及其返回结果进行染色与时间轴排列，可逐条展开查看。

为支撑下钻视图对推理轨迹的稳定关联，模块在任务执行前后对智能体会话目录做快照对比，把本次任务新写入的会话日志路径记录到任务元数据中，从而避免依赖文件名匹配这种脆弱方式建立任务、流量条目与推理轨迹之间的对应关系；同时为防止单条轨迹文件过大导致前端内存压力，加载阶段设有字节阈值并对超长记录进行截断保护，超出部分以截断标志而非静默丢弃的方式呈现给审计人员。借助上述机制，可解释性不再依赖事后构建的解释模型，而是直接以推理过程的原生产物形式被纳入审计视图。

模块在职责划分上严格遵循“合并层只对齐不判别，统计层只派生不修改，呈现层只渲染不二次推断”的边界约定。所有上游产物均以只读方式读取，本模块不会对预处理或Agent的中间结果进行任何二次推断，从而保证审计结果与中间证据之间的因果链条清晰、可复核。

\section{系统工程实现}[System Engineering Implementation]

本节从工程实现视角介绍系统的目录组织、配置管理、端到端运行流程、异常处理策略与测试机制，说明上述模块如何在统一的项目结构下协同工作。

项目根目录采用按职责分层的方式组织。核心代码包括预处理、合并、可视化三个顶层组件，以及Agent适配、模型适配与通用工具三个子包，分别承担规则化筛选、Agent适配、模型适配与底层辅助职能；分类模型作为自包含子项目，封装了推理脚本、模型结构、权重文件与词表；智能体引擎以可安装包形式接入，由业务侧通过适配层调用；智能体工作区作为隔离运行目录，包含角色契约、技能定义、配置文件以及会话与输入输出存储；脚本目录提供端到端流水线脚本与各模块单独验证脚本；数据目录集中管理样例PCAP、Zeek/nDPI中间产物、未知流PCAP切片以及最终报告；测试目录存放面向合并模块与解析器的单元测试。这一目录结构使得项目在保持核心代码集中的同时，能够清晰地隔离子项目、运行时数据与测试脚本，便于定位与维护。

配置管理集中在统一的配置入口中，遵循“代码内只定义路径，敏感参数走环境变量”的原则。配置入口首先从环境变量读取Zeek与nDPI可执行文件路径，再读取分类模型相关的启用标志、适配器路径、权重路径与词表路径；接着定义由项目根目录衍生出的全部数据目录，并在导入时确保所有目录就绪；最后明确智能体工作区路径与引擎配置路径，避免使用用户级全局配置带来的隔离性问题。配置入口还提供布尔值与JSON结构的环境变量解析辅助函数，统一处理跨环境部署中的差异，提升了系统的鲁棒性。

端到端运行由专门的流水线驱动器串联完成，整体流程见算法\ref{alg:ch4-pipeline}。流水线接收PCAP路径、是否跳过Agent或可视化、Agent单批最大流数等命令行参数，按四个阶段顺序执行；并在控制台同时打印任务概要，包括已知/未知/Agent标注数量、最终报告路径以及各图表文件位置，方便运维人员核查。

  \begin{algorithm}[h]
	\caption{端到端流水线执行流程}\label{alg:ch4-pipeline}
	\KwIn{PCAP 文件路径 $p$, 运行参数 $\Theta$}
	\KwOut{最终审计报告 $R$}

	初始化预处理器并执行 $P \gets \mathrm{Preprocess}(p)$\;
	切分未知流为单流 PCAP 并写入任务清单\;
	\If{分类模型已启用}{
		调用适配函数获得分类结果 $C$ 并回填到 $P$\;
	}
	\eIf{Agent 阶段未跳过 且 $P$ 包含未知流}{
		对未知流执行可识别性预筛分，得到可推断子集 $U_a$ 与无信号子集 $U_d$\;
		对 $U_d$ 直接生成附带具体缺失原因的空决策\;
		将 $U_a$ 切分为多个分块并以受限并发的异步方式调用智能体\;
		对智能体输出执行宽容解析与流键白名单过滤\;
		对最终标签为空且分类模型概率达到阈值的决策按式\eqref{eq:ch4-fallback}合成兜底标签，得到 $A$\;
	}{
		$A \gets \varnothing$\;
	}
	调用合并函数生成最终报告 $R \gets \mathrm{Merge}(P, A)$\;
	\If{可视化阶段未跳过}{
		调用绘图接口生成饼图、柱状图、桑基图并回写到 $R$\;
	}
	持久化 $R$ 与中间产物，输出任务概要\;
	\Return{$R$}\;
\end{algorithm}

异常处理在系统中按层次组织。最底层的Zeek与nDPI调用通过子进程接口捕获返回码与标准错误，对超时、可执行文件缺失、CSV格式异常等情况显式抛出运行时异常并附带可读信息；中间层的分类模型适配器通过异常拦截机制将子进程错误转化为带有错误标识的占位结果，使主流程能够继续；智能体运行驱动器为每个分块单独捕获异常，将错误以空决策数组形式写入输出文件，避免某个分块失败导致整体流水线中断；最上层的流水线驱动器将所有阶段性异常收集到错误列表，最终一并写入最终报告的错误字段，使审计人员能够在结果中直接看到失败原因，而不必反查日志。系统对外的所有错误信息均使用结构化字段表达，便于后续做自动化告警或质量追踪。

测试与脚本验证由两层组成。单元测试层面使用Pydantic模式校验与字典断言为合并函数与结果解析器提供测试用例，覆盖标签拆分、置信度截断、服务类型词表校验、聚合排序等关键行为；模块级脚本层面则在真实样例PCAP上对各模块独立验证，包括Zeek与nDPI调用、预处理结果生成、Agent连通性与端到端任务执行等场景。其中Agent验证脚本提供连通性检查、给定输入分块的端到端验证以及自定义提示词调试三种工作模式，分别用于服务可达性、端到端正确性与提示词迭代调试。这些脚本与单元测试形成互补，前者保证逻辑正确性，后者保证模块可在真实数据上稳定运行。

通过上述目录组织、配置体系、流水线驱动、分级异常处理与多层测试机制，系统在工程层面具备了模块化、可调试、可降级、可复现的运行能力，能够在不同实验环境与部署条件下输出一致的分析结果。

\section{本章小结}[Chapter Summary]

本章围绕“面向大规模网络流量的智能应用分析系统”的设计与实现展开了系统化论述。首先在总体架构层面给出了“两阶段混合\,+\,插件式扩展”的设计思想，明确了已知流量过滤、流量分类、智能分析、结果合并与可视化等核心模块的职责边界以及彼此之间的数据契约。接着分模块说明了关键实现细节：已知流量过滤模块以Zeek与nDPI为基础协议解析层，结合双键索引、五元组流聚合与多级判别规则完成确定性筛选，并产出可被分类模型直接消费的单流PCAP与任务清单；流量分类模型模块以适配函数协议进行可插拔接入，依托预训练Transformer编码器对未知流量给出五类服务类型概率证据，并通过两轮迭代式微调缓解了域外分布漂移问题；智能分析模块在隔离工作区与契约化输入输出的基础上，通过可识别性预筛分降低无效推理、通过工具路由曝光与运行时调用预算约束工具使用质量、通过分块级异步并发提升吞吐、通过宽容解析与流键白名单提升输出鲁棒性，并以分类模型概率作为大语言模型空标签时的兜底证据，实现了语义判断与统计判断的分层融合；结果合并与可视化模块以流唯一键聚合两阶段产物，输出聚合统计与饼图、柱状图、桑基图三类可视化结果，并将智能体推理轨迹作为可追溯的解释来源纳入审计视图。最后从工程实现角度阐述了项目目录、配置体系、端到端流水线、分级异常处理、推理轨迹索引与多层测试机制等支撑机制，验证了系统在模块化、可扩展性、可解释性与容错性方面对第3章需求分析的对应。整体设计为后续实验评估与实际部署提供了完整的工程基础。